% ====================================================================
% paper_body.tex -- specialist-journal article (PRD working target).
% Owner: publication lane (docs/publication/). The s5p campaign owns the
% Stage-7 inference text in the note and primer and sends paper-side
% wording here (agreed 2026-10-06). Claim-to-evidence routes:
% docs/publication/CLAIMS-20261005-claim-to-evidence.md.
% The full technical and diagnostic record lives in the analysis note.
% ====================================================================

\section{Introduction}

Differential neutrino cross sections require correcting observed event samples
for detector response.  Iterative Bayesian unfolding on a binned migration
matrix is widely used, but the response tensor grows rapidly when additional
observables are reported.  OmniFold instead learns event weights in an unbinned
feature space through classifier-based density-ratio estimation
\cite{Andreassen:2019cjw,Andreassen:2021zzk}.  This makes simultaneous
higher-dimensional unfolding possible without specifying a separate binned
response tensor for every projection.  Unbinned unfolding has been applied to
collider data, for example eight observables in deep-inelastic scattering at
H1~\cite{H1:2021wkz} and twenty-four $Z$+jets observables at
ATLAS~\cite{ATLAS:2024jets}; see Ref.~\cite{Canelli:2025guide} for a survey.  In
neutrino physics it has been demonstrated on a public simulated T2K near-detector
sample~\cite{Huang:2025omnifold}.  We are not aware of an application to
neutrino-scattering data.

We apply the method to MINERvA medium-energy forward-horn-current $\nu_{\mu}$
data from the MINERvA open-data release~\cite{MINERvA:OpenData}.  We first
reproduce the published double-differential charged-current inclusive cross
section in transverse and longitudinal muon momentum,
$d^{2}\sigma/(d\pT\,d\pPar)$~\cite{MINERvA:2021owq}, then add available
hadronic energy $\Eavail$, three-momentum transfer $q_{3}$, and hadronic
invariant mass $W$ to obtain simultaneous three-, four-, and five-observable
central values.  Multi-differential neutrino measurements on data reach three
dimensions~\cite{MINERvA:2022qelike3d,MINERvA:2026qeleme,MicroBooNE:2023inc3d,NOvA:2026inc3d}.
The five observables here are not five independent measurements: $\Eavail$,
$q_{3}$ and $W$ are all built from reconstructed hadronic energy.

The paper has two parts.  The first validates the unbinned construction against
the published two-dimensional result and characterizes its higher-dimensional
central values, including measured limitations: a fixed-truth undercoverage of
the two-dimensional statistical band as first built, a background-subtraction bias at high $W$,
and a regularization bias under generator-sized shape changes.  Because the
estimator's model dependence prevents a qualified total uncertainty at useful
precision, we do not report a five-dimensional cross-section measurement with
uncertainties.  The second part instead tests specified generator predictions
directly: a joint statistic on the five-dimensional result is calibrated by
pseudo-experiments that run the full analysis chain under a declared nuisance
model (Sec.~\ref{sec:joint}).

\section{Data, signal and simulation}

The analysis uses the twelve MINERvA medium-energy playlists, corresponding to
$1.057\times10^{21}$ protons on target and an average neutrino energy near
\SI{6}{GeV}, from the MINERvA open-data release of pre-selected
AnaTuples~\cite{MINERvA:OpenData}; the detector is described in
Ref.~\cite{MINERvA:2013zvz}.  The files analysed are an earlier production of
that release; their identities are recorded for reproduction
(Sec.~\ref{sec:data-availability}).  The signal definition, fiducial volume, flux
treatment~\cite{MINERvA:2016iqn}, and reconstruction-level selection follow
Ref.~\cite{MINERvA:2021owq}.  After selection, the sample contains $4.12$
million data events, and the simulation contains $32.8$ million true signal events, generated with
GENIE~2.12.6 and the MINERvA Tune~v1, which reweights GENIE with a
non-resonant pion reduction, the RPA correction and an empirical low-recoil
two-particle--two-hole enhancement~\cite{MINERvA:2016lowrecoil}.

Reconstructed $\Eavail$ is the calorimetric recoil energy in the tracker and
electromagnetic calorimeter after removal of the muon's energy deposits; the
reconstructed energy transfer $q_{0}$ is the calorimetric recoil energy, from
which $Q^{2}$, $q_{3}$ and $W$ follow with the muon kinematics.  The external
generator predictions compared below are GENIE~2.12.10 on CH without and with
the Valencia two-particle--two-hole component, NuWro~21.09 and GiBUU~2019 on
carbon.  The GENIE and NuWro predictions use a repaired flux sampling; the GiBUU
sample stops at $E_{\nu}\approx\SI{20}{GeV}$.

\section{Unbinned unfolding}
\label{sec:method}

Each OmniFold iteration consists of two classification steps.  A
reconstruction-level classifier distinguishes data from simulated events and
defines a likelihood-ratio weight.  The event pairing transfers that weight to
truth, including unmatched events, and a truth-level classifier propagates a
smooth reweighting to the next iteration.  Cross sections are formed by
histogramming the reweighted truth events in a declared reporting binning and
dividing by integrated flux, target count, and bin width.

Four estimator configurations appear in this paper, and each result names the
one it uses.
\begin{itemize}
\item \emph{Two-dimensional central value} (Sec.~\ref{sec:validation}): five
iterations of depth-three gradient-boosted decision trees with 100 trees, fitted
with exact-split scikit-learn gradient boosting, and a purity-based background
subtraction.
\item \emph{Three- to five-dimensional central values}
(Secs.~\ref{sec:validation}--\ref{sec:central}, Figs.~\ref{fig:joint}
and~\ref{fig:generator-context}): the same five-iteration configuration fitted
with LightGBM at matched nominal capacity, with the purity-based background
subtraction.
\item \emph{Two-dimensional uncertainty ensembles}: the finalized systematic-universe
and Poisson-bootstrap ensembles use LightGBM at matched nominal capacity, so
the two-dimensional central value and its covariance are not estimator-matched.
\item \emph{Joint-test estimator} (Sec.~\ref{sec:joint}): a deterministic
LightGBM OmniFold, five iterations, with an unbinned negative-weight background
treatment whose refinement classifier uses 400 trees and 31 leaves.  The
negative-weight treatment reduces the purity method's nominal-truth bias in the
highest-$W$ cells about fifteen-fold, and the larger refinement classifier
removes most of the background-related bias that remains; the configuration was
selected before the joint-test design was frozen.
\end{itemize}

\section{Validation}
\label{sec:validation}

\subsection{Two-dimensional reproduction}

The two-dimensional result provides a stringent end-to-end validation.  Its
integrated cross section is \SI{\sigTwoD}{cm^2/nucleon}, compared with
\SI{\sigTwoDpaper}{cm^2/nucleon} in the published analysis, a ratio of
$\ratioTot$.  Of the 205 reported bins, \SI{\binsTen}{\percent} agree within
\SI{10}{\percent}; the standardized bin differences formed with the published
uncertainties have an RMS of $\pullRMS$; against the full published covariance
the indicative distance is $\chi^{2}/\mathrm{ndf}=\chiPaper$, a correlated
shape difference along the low-$\pT$ peak ridge.  Figure~\ref{fig:validation}
shows the principal one-dimensional projections.  Because both results derive
from the same data, this validates the estimator rather than comparing
independent measurements; no combined-covariance statistic is quoted.

\begin{figure*}[t]
  \centering
  \includegraphics[width=0.78\textwidth]{model_comp_projections}\\[-0.4em]
  \begin{minipage}[c]{0.31\textwidth}
    \centering
    \includegraphics[width=\linewidth]{paper_validation_residual}
  \end{minipage}\hspace{1.2em}
  \begin{minipage}[c]{0.30\textwidth}
    \centering\footnotesize
    \textbf{Median relative uncertainty}\\[0.35em]
    \begin{tabular}{@{}lr@{}}
      This work & \SI{\uqMedian}{\percent} \\
      Published & \SI{\uqPaper}{\percent}
    \end{tabular}\\[0.8em]
    \textbf{Published-$\sigma$ standardized residuals}\\[0.35em]
    mean $=\pullMean$, RMS $=\pullRMS$
  \end{minipage}
  \caption{End-to-end validation against the published binned result
  \cite{MINERvA:2021owq}.  Top: $\pT$ and $\pPar$ projections of the published
  result, the five-iteration OmniFold result, and the MINERvA Tune~v1.  The two
  measurements use the same data, so their overlay validates the estimator
  rather than comparing independent results.
  Bottom left: $(\mathrm{OmniFold}-\mathrm{published})/\sigma_{\mathrm{published}}$
  on the reported two-dimensional grid; it is a descriptive shared-data
  residual, not an independent-result pull.  Bottom right: the independently
  reconstructed and published median relative uncertainties.}
  \label{fig:validation}
\end{figure*}

\subsection{Two-dimensional uncertainty and its coverage}

The standalone covariance combines matched-central-value systematic universes
with statistical and training-seed blocks.  Its median relative uncertainty is
\SI{\uqMedian}{\percent}, compared with \SI{\uqPaper}{\percent} for the
publication.  The coverage of this total is untested.  Its statistical block,
as first built, was tested against a fixed truth in \covTwoDNtoy{} pre-registered
closure pseudo-experiments and undercovered in its tails:
\SI{\covTwoDCTwoPctOne}{\percent} of standardized residuals fell within two standard
deviations, against \SI{95.4}{\percent} nominal, while the one-standard-deviation
fraction, $\covTwoDCOne$, was nominal; the pull mean was consistent with zero.  The
main cause, a completeness division that the bootstrap replicas carried and the
central value does not, was found and removed.  The rebuilt statistical band, now
the quoted one, has a median statistical uncertainty of \SI{0.674}{\percent} and
leaves the total median at \SI{\uqMedian}{\percent}.
Rescoring the same closure pseudo-experiments against the rebuilt band, which is
post hoc and partly circular, removes the tail undercoverage
(\SI{97.7}{\percent} within two standard deviations), while \SI{78.1}{\percent} fall
within one standard deviation, above its nominal window, so the rebuilt band is
wider than the toy scatter; the residual excess width lies in the data-statistics
component and is not explained by background purity, and its cause is untested.
The rebuilt band has not been independently re-tested.
% Source: docs/orchestration/OUTCOME-20261005-2d-fixed-truth-coverage-fail.md;
% VALIDATION_LEDGER VL169 (test of the earlier band); VL170 (rebuilt band, adopted by Joseph
% 2026-10-06, docs/publication/DECISION-20261006 addendum af24a101); VL172 (rebuilt products;
% main 2c9939fe; the quoted 6.87% chain is 2d-unfolding/uq/rollup_vl170_adoption.sh /
% recompute_2d_budget.py); rescoring 5aab2d6d (descriptive, post hoc); purity check receipt
% purity_datastream_check.json; KNOWN_ISSUES 84, 85.

\subsection{Higher-dimensional central values}

Every higher-dimensional central-value construction is tested against its
lower-dimensional predecessor.  The three-, four-, and five-observable marginal
normalizations agree at the approximately $1$--$2\%$ level, and a closure test
recovers an injected $\Eavail$ distortion.  Such closure tests reweight the
simulated signal and do not exercise the background subtraction.  End-to-end
pseudo-experiments that include the purity-based background subtraction
reproduce the input at nominal truth except in the highest-$W$ cells, which are
biased by up to \SI{4}{\percent} in simulation, about one quoted standard
deviation in one $(\Eavail,W)$ cell; this bias is not included in the quoted
uncertainties.  The unbinned negative-weight background treatment, as first
measured with a deterministic variant of the five-dimensional LightGBM production estimator, reduces it to
below \SI{0.3}{\percent}, leaving up to \SI{1}{\percent} in some joint cells; on
data the two treatments differ by at most \SI{1.4}{\percent}.
% Source: OUTCOME-20260925-s5c-purity-background-bias-at-high-W.md; VL149,
% VL151, VL152; KNOWN_ISSUES 75.

The same pseudo-experiments show that the unfolding follows a generator-sized
change of the $\Eavail$ shape only partly.  When the true shape is reweighted
by the GiBUU-to-GENIE ratio (a test deformation of that size rather than a
measured generator difference: the GENIE prediction it used was later found to
have been generated from a mis-sampled flux), the $(\Eavail,W)$ cells deviate
from it by several percent in the largest cells and by up to \SI{74}{\percent}
at high $W$, with or without background.  This regularization bias is not in
the quoted uncertainties.  The tested iteration and capacity scans do not
remove it; for other fixed generator shapes it is \SIrange{16}{31}{\percent} in
the most affected cell and within the quoted total uncertainty in about half of
the $(\Eavail,W)$ cells.  The unfolded values on data also move by up to a few
statistical standard deviations under input changes at the level of
floating-point rounding; a later measurement found this spread already inside
the tested candidate's data-bootstrap statistical uncertainty, which refits
every replica, and well below the total uncertainty, but it limits how exactly
the central values can be reproduced.  No frequentist coverage is claimed for
the five-dimensional intervals: an end-to-end test of bias-corrected
statistical intervals from a deterministic variant of the estimator failed its
predeclared criterion at nominal truth and under an $\Eavail$ tilt.
% Sources: OUTCOME-20260925-s5n-stage1-development-fail.md (VL151, VL152);
% OUTCOME-20260926-s5e-oi192-diagnosis-and-candidate.md (VL153, VL154);
% RECORD-20260927-s5p-stage2-exit.md (VL155); OUTCOME-20260925-s5c-tier-s-futility-fail.md (VL150).

\section{Central-value results}
\label{sec:central}

The $\Eavail$ dimension provides an independent view of the established
low-recoil discrepancy~\cite{MINERvA:2016lowrecoil,MINERvA:2022incl}.  GENIE,
the MINERvA Tune~v1, NuWro, and GiBUU all underpredict the unfolded central
value at low $\Eavail$, between the quasielastic and $\Delta$ peaks.  Pion
final-state-interaction variations (same-sample ratios on pre-repair GENIE)
change $d\sigma/d\Eavail$ by less than one percent, whereas the empirical
Valencia two-particle--two-hole contribution fills \SI{52}{\percent} of the
data--GENIE gap below $\Eavail=\SI{0.4}{GeV}$.  These tests constrain the chosen
pion-FSI variations, not all nuclear-transport models.  Adding $q_{3}$ resolves
the momentum-transfer dependence of this feature: on the maximal common
$(\Eavail,q_{3})$ grid, the central value lies within \SI{10}{\percent} of the
dedicated published low-recoil measurement~\cite{MINERvA:2022incl}.
% Sources: VL159, VL160.

The fifth observable, $W$, separates the remaining discrepancy from the
low-mass enhancement.  Figure~\ref{fig:joint} shows the joint $(\Eavail,W)$
projection relative to the MINERvA Tune~v1.  In cell-integrated terms, two
thirds of the positive difference lies at $\Eavail\ge\SI{0.8}{GeV}$ and most of
that at $W\ge\SI{1.8}{GeV}$, but a fifth of it sits in the single widest catch
cell; in ratio, the data exceed Tune~v1 by \SIrange{12}{30}{\percent} in those
cells and by \SIrange{23}{31}{\percent} at $W<\SI{1.1}{GeV}$ for every $\Eavail$,
so the relative excess is not confined to high $W$.  These statements are
descriptive: no significance is assigned, no response-mismatch closure has been run
for this map, and the high-$\Eavail$, high-$W$ region is where the regularization
bias of Sec.~\ref{sec:validation} is largest.
% Source: shares 67%/83% (sec_eavailw.tex, cell-integrated) and the catch-cell 22% and ratio ranges
% recomputed from the released figure arrays (fig_arrays.npz sha256 72394a2b..., release/article-figures-20261006).

\begin{figure}[t]
  \centering
  \includegraphics[width=\columnwidth]{paper_joint_localization}
  \caption{Joint $(\Eavail,W)$ central value (five-dimensional LightGBM production estimator), on the
  $7\times6$ reporting grid with $\Eavail$ edges
  $[0,0.1,0.2,0.4,0.8,1.5,3.0,100]\,\si{GeV}$ and $W$ edges
  $[0,1.1,1.4,1.8,2.2,3.0,100]\,\si{GeV}$; axes are cell indices counting from
  zero, so index~6 in $\Eavail$ and index~5 in $W$ are wide catch bins.  Left:
  the unfolded-to-Tune~v1 ratio per cell, dimensionless, on a $0.5$--$1.5$
  scale.  Right: the unfolded-minus-Tune~v1 difference as a
  \emph{cell-integrated} cross section in \si{cm^2/nucleon}, i.e.\ the
  density difference multiplied by that cell's
  $\Delta\Eavail\,\Delta W$; the catch-bin cells therefore carry their large
  widths and appear larger here than the ratio panel alone would suggest.  No
  significance is assigned.  The comparator is the MINERvA Tune~v1 (GENIE~2.12.6
  with the analysis tune weights), not the untuned GENIE central value of
  Fig.~\ref{fig:generator-context}.}
  \label{fig:joint}
\end{figure}

The one-dimensional projections in Fig.~\ref{fig:generator-context} compare the
same central result with four external predictions: GENIE's own central value,
the same GENIE with the Valencia two-particle--two-hole component, NuWro, and
GiBUU.  All four lie below the data overall (integrated data-to-prediction
ratios $1.07$--$1.39$) and at high $W$.  For GENIE and NuWro the high-$W$
shortfall is about their overall one (ratios $1.11$--$1.13$ for
$2.2\le W<\SI{3.0}{GeV}$, against $1.07$--$1.15$ integrated), and in the joint
high-$\Eavail$, high-$W$ region ($\Eavail\ge\SI{0.8}{GeV}$, $W\ge\SI{1.8}{GeV}$)
it is $1.14$--$1.16$, within \SI{7}{\percent} of their integrated ratios.
GiBUU, the most deficient, falls furthest behind there ($1.61$, against $1.39$
overall), but its sample stops at $E_\nu\approx\SI{20}{GeV}$: in GENIE and NuWro
higher energies carry \SIrange{14}{15}{\percent} of this region's cross section
and \SI{7}{\percent} of the total, and supplying that share gives GiBUU about the
others' contrast ($1.37$ against $1.29$).  Turning on the Valencia
two-particle--two-hole component leaves that region unchanged because it
primarily populates lower $W$.  The high-$\Eavail$, high-$W$ difference of Fig.~\ref{fig:joint} is thus
a statement relative to the MINERvA Tune~v1.  Dedicated MINERvA
shallow-inelastic measurements also report generator
disagreements~\cite{MINERvA:2026sis}; their different selections preclude
treating them as independent confirmation of this high-$\Eavail$, high-$W$ difference.
% Sources: VL156-VL161; KNOWN_ISSUES 82.

\begin{figure*}[t]
  \centering
  \includegraphics[width=0.84\textwidth]{paper_eavailW_generators_ratio}
  \caption{Generator context for the unfolded $\Eavail$ (left) and $W$ (right)
  central-value projections (five-dimensional LightGBM production estimator), in
  \si{cm^2/GeV/nucleon}, with the data-to-prediction ratio below each; the wide catch
  bins ($\Eavail>\SI{3}{GeV}$, $W>\SI{3}{GeV}$) are omitted from the display.  Drawn from
  the released figure arrays (Sec.~\ref{sec:data-availability}).  GENIE-CV (untuned
  GENIE~2.12.10 on CH, with no two-particle--two-hole events), the same GENIE
  with the Valencia two-particle--two-hole contribution, NuWro~21.09, and
  GiBUU~2019 (both on carbon; GENIE and NuWro flux-repaired) all remain below the
  unfolded result overall and at high $W$ (the two GENIE predictions lie slightly above it
  at $1.4\le W<\SI{1.8}{GeV}$); GENIE and NuWro by about their overall
  normalization shortfall, GiBUU, whose sample stops at $E_\nu\approx\SI{20}{GeV}$,
  by more.  No significance is assigned.}
  \label{fig:generator-context}
\end{figure*}

\section{Calibrated joint tests of generator predictions}
\label{sec:joint}

\subsection{Hypotheses, statistics and calibration}

For each of five predictions $G$ --- the MINERvA Tune~v1, GENIE~2.12.10 without
and with the Valencia two-particle--two-hole component, NuWro~21.09 and
GiBUU~2019 --- the null hypothesis $H_{0}(G)$ is that the true cross section equals
$G$ on a fine five-dimensional grid, with the MINERvA Tune~v1's shape and $G$'s
finite-sample fluctuations below that grid.  For the Tune~v1, the analysis
simulation itself, this is $H_{0}$ exactly.  The tests are evaluated on a coarse
$3^{5}$ grid in $(\pT,\pPar,\Eavail,q_{3},W)$ with edges
$\pT\in[0,0.55,0.85,4.5]\,\si{GeV/c}$, $\pPar\in[1.5,3.5,6,60]\,\si{GeV/c}$,
$\Eavail\in[0,0.4,1.5,100]\,\si{GeV}$, $q_{3}\in[0,1.2,2,100]\,\si{GeV}$ and
$W\in[0,1.4,2.2,100]\,\si{GeV}$, restricted to the \jtNcells{} supported cells;
for GiBUU, whose sample lacks high-energy neutrinos, the domain is the
\jtNcellsGiBUU{} cells with $\pPar<\SI{6}{GeV/c}$.
% Source: amendment 7 claims.family; nd-unfolding/s5p_stage1_inspect.py:45 (J_EDGES).

Two statistics are formed from the unfolded cell values $F$ and the prediction
$\mu_{G}$: a total statistic $T_{\mathrm{tot}}=r^{\top}W^{-1}r$ with
$r=F-\mu_{G}$, and a shape statistic $T_{\mathrm{sh}}$ of the same form on the
normalized cell fractions, with the normalization Jacobian applied to $W$ and one
cell dropped.  The metric $W=V+\mathrm{diag}\,\mathrm{Var}(\mu_{G})$ adds the
prediction's own finite-sample variance to $V$, a Ledoit--Wolf-shrunk
covariance~\cite{LedoitWolf:2004} of \jtVn{} pilot pseudo-experiments at the
Tune~v1 null.  $V$ sets the power of the test, not its size; it is not a
measurement covariance.
% Source: nd-unfolding/s5p_inference.py:30-49; nd-unfolding/s5p_joint.py:147-162;
% amendment 7 calibration.metric and metric_frozen.

The null distribution of each statistic is obtained from pseudo-experiments that
run the full analysis chain at the null truth.  Each draws Poisson pseudo-data
from one half of the simulation together with an independent background source,
refits the background-refinement classifier, and unfolds with the joint-test
estimator trained on the other half.  Every experiment also draws one of
\jtFluxUniverses{} flux universes, every one of the \jtModelBands{}
interaction-model uncertainty bands at one of its universes, the weight-based
detector bands (MINOS efficiency and GEANT neutron, pion and proton interaction
reweights), and the five muon and beam-angle lateral bands through a linear
surrogate, together with a \SI{1.4}{\percent} normalization draw, the data's rounding
noise taken from 20 real-data re-unfolds with rounding-level input changes, and the
predictions' finite-sample residual.  The data are unfolded once with the full
simulation.  The Monte Carlo
$p$ value of a test is $(k+1)/(B+1)$, where $k$ of $B$ null pseudo-experiments
reach the observed statistic; it is never zero.  The claim $p$ value is the
largest over declared variants that move the null by an upper bound on the
simulation-size difference between the pseudo-experiments and the data unfold
and by $\pm2$ times an estimate of the residual below the fine grid.  The
external nulls are strongly non-central, with non-centralities $\lambda$ of
\jtLambdaMin{}--\jtLambdaMax{} in these units, because the calibration
reproduces the estimator's pull toward its prior.
% Source: amendment 7 calibration.process, calibration.process_shift, m1_shift,
% prefreeze_measurements.summary.

The ten tests (two per prediction) are controlled at a familywise level
$\alpha=\jtAlpha$ by the Holm procedure~\cite{Holm:1979}, with a determinacy
requirement: a step rejects only when the 95\% Clopper--Pearson
interval~\cite{Clopper:1934} of $k/B$ lies entirely below its Holm threshold.
Calibration proceeded in batches of \jtBatch{} pseudo-experiments per null, to
at most \jtBcap{}, under a sequential rule that stopped a null once both of its
decisions were determined.  The hypotheses, statistics, variants, multiplicity
rule, stopping rule and every input were frozen before any production
pseudo-experiment or observed statistic.  Two data-informed elements are
disclosed: the design reviewers had seen descriptive per-cell residuals before
the freeze, and the controller's status files exposed observed claim $p$ values
during production; a bounded extension of the production budget was decided
while they were visible, to let the sequential rule reach its frozen stop rather
than a budget stop.  No statistic, variant, threshold or stopping rule was changed,
and the Monte Carlo $p$ values remain valid at the number of pseudo-experiments
reached.
% Source: amendment 7 data_informed_disclosure and calibration.sequential_rule;
% DECISION-20261004-s5p-production-budget-extension.md (Joseph, 2026-10-04: "Keep all scientific
% criteria unchanged").  An independent recomputation of every
$p$ value and decision from the stored products, with separately written code,
agrees with the production evaluation in all \jtRecompRows{} compared quantities.
% Sources: amendment 7 claims.rejection, calibration.sequential_rule,
% validation_and_assurance (v); REPORT-20261005-s5p-recompute-final-verification.md sec. 8.

\subsection{Conditions}

These are statements about the simple hypotheses $H_{0}(G)$, not a measured
cross section, and no new reportable uncertainty is produced.  Each decision
below holds under the conditions of the declared model, stated here because they
bound what the tests establish.  (i) Below the fine grid the
null truth carries the Tune~v1's shapes; the fine-grid convergence check is not
converged (successive refinements change the result by
\jtMoneMin--\jtMoneMax{} of the previous change), and its effect enters as the
residual variants above.  (ii) The detector model is the simulation's, varied
only by the drawn detector bands; the analysis has no band for the
calorimetric recoil-energy response, so completeness of the hadronic-response
model is not established (Sec.~\ref{sec:response}).  (iii) The interaction-model
prior is the analysis's 34-band set, with negative weights of two bands clipped per
universe.  (iv) The pseudo-experiments use half of the simulation and the data
unfold all of it; the resulting shift is bounded by a declared variant.  (v)
GiBUU lacks $E_{\nu}>\SI{20}{GeV}$, and its out-of-domain truth enters the
migrations into its domain.  For all predictions $E_{\nu}>\SI{100}{GeV}$ is absent, which is negligible.  (vi) The predictions' own finite samples enter the
metric and the simulated residuals.
% Source: amendment 7 claims.conditions_stated_with_every_claim (all six).

\subsection{Results}

Table~\ref{tab:joint} lists the ten decisions as recorded.  Every test is
rejected at familywise $\alpha=\jtAlpha$; nine have $k=0$, and NuWro's shape test
has $k=1$.  Each rejection survives when the sub-fine-grid residual variants are
enlarged from two to three times their estimate.  Figure~\ref{fig:joint-nulls} shows each null's
calibrated distribution with the observed statistic.  The Tune~v1 null, the analysis
simulation itself, is centred at small values, whereas the external nulls are
displaced far upward by the estimator's pull toward its prior, which the
calibration reproduces; the comparison of the data with each null is made against
that displaced distribution.

\begin{figure*}[t]
  \centering
  \includegraphics[width=0.98\textwidth]{paper_joint_nulls}
  \caption{Calibrated null distributions of the joint statistics (joint-test
  estimator) for each prediction (columns), total statistic (top) and shape
  statistic (bottom), shown for the declared claim variant whose null distribution
  has the largest median; the vertical line is the observed statistic and $k/B$ is
  as in Table~\ref{tab:joint}.  The horizontal axis is logarithmic.  Plotted from
  the released sufficient inputs (Sec.~\ref{sec:data-availability}), which reproduce
  every recorded $p$ value and decision.}
  \label{fig:joint-nulls}
\end{figure*}

\begin{table}[t]
  \caption{Joint-test decisions as recorded (Tune~v1 null: the analysis
  simulation; $B$: null pseudo-experiments completed at the frozen stop, of
  1400 submitted per null (1800 for NuWro); $p=(k+1)/(B+1)$;
  CP: 95\% Clopper--Pearson interval of $k/B$; all ten are robust to the
  sub-fine residual at three times its estimate).  Every submitted
  pseudo-experiment, including those first lost to time limits, has been observed;
  none of the recovered ones reaches the observed statistic (Sec.~\ref{sec:missing}).}
  \label{tab:joint}
  \centering\scriptsize
  \begin{tabular}{@{}llrrlll@{}}
    \toprule
    $G$ & test & $B$ & $k$ & $p$ & Holm thr. & CP interval \\
    \midrule
    NuWro 21.09 & total & 1751 & 0 & 1/1752 & 0.005 & $[0, 0.00210]$ \\
    GENIE CV & total & 1366 & 0 & 1/1367 & 0.00556 & $[0, 0.00270]$ \\
    GENIE CV & shape & 1366 & 0 & 1/1367 & 0.00625 & $[0, 0.00270]$ \\
    Tune~v1 & total & 1365 & 0 & 1/1366 & 0.00714 & $[0, 0.00270]$ \\
    Tune~v1 & shape & 1365 & 0 & 1/1366 & 0.00833 & $[0, 0.00270]$ \\
    GiBUU 2019 & total & 1351 & 0 & 1/1352 & 0.01 & $[0, 0.00273]$ \\
    GiBUU 2019 & shape & 1351 & 0 & 1/1352 & 0.0125 & $[0, 0.00273]$ \\
    GENIE + MEC & total & 1343 & 0 & 1/1344 & 0.0167 & $[0, 0.00274]$ \\
    GENIE + MEC & shape & 1343 & 0 & 1/1344 & 0.025 & $[0, 0.00274]$ \\
    NuWro 21.09 & shape & 1751 & 1 & 2/1752 & 0.05 & $[0.0000145, 0.00318]$ \\
    \bottomrule
  \end{tabular}
  % Source: RECORD-20261005-s5p-joint-5d-inference-result.md sec. 2 (origin/main 9b26b8c3),
  % joint-evaluate.json sha256 b9604502...; robustness: robust-labels.json 206655f9...
\end{table}

The published two-dimensional comparison on the same data already disfavours
these generator families strongly~\cite{MINERvA:2021owq}, so the rejections
are not by themselves a new statement that they fail.  What the joint tests
add is a calibrated statement on the five-observable result that accounts for
the estimator's own bias toward its prior and for the declared nuisances.
Restricting the same calibrated tests to matched $3\times3$ coarse projections of
the same cells, onto $(\pT,\pPar)$ and onto $(\Eavail,W)$, gives $p<0.05$ for
every test in at least one of the two projections, so we make no claim that the
joint five-observable information adds discrimination beyond them.  This comparison
was specified before it was run but after the five-dimensional outcomes were known;
it uses the same calibration ensembles, variants and claim rule, and its
$p$ values were reproduced independently.
% Source: docs/publication/w1/W1-RECORD-20261006.md (22a004e7, reproduction 09718448).
What the rejections establish is therefore narrow.  They are global statements
about five specified predictions under the declared model; they do not identify a
kinematic region or a nuclear mechanism, they say nothing about more recent
generator versions, and each rejected test already has $p<0.05$ in at least one
of the two matched coarse projections.  Power is context only, because every test rejects: at the Tune~v1
null, two of three declared alternatives are detected in shape with power 1.0 at the
0.005 level, and the third, a MEC-like $(\Eavail,W)$ shape, with 0.23 (each power set
lost some of its 200 alternatives to time limits, retaining 172 to 199); at the
GENIE~2.12.10 null the shape powers are 0.74, 0.02 and 0.00.
% Source: s5p Stage-7 approved wording, sec_eavailw.tex sec. 8.6 (sec:joint5d) at main cbd075e7;
% REVIEW-20261006-s5p-stage7-note-wording.md (two rounds, findings resolved).

\subsection{Lost pseudo-experiments}
\label{sec:missing}

Scheduler time limits interrupted or prevented \jtLostCal{} of the submitted
calibration pseudo-experiments (\jtLostInterrupted{} interrupted,
\jtLostNever{} never started) and \jtLostPow{} power pseudo-experiments, so the
decisions in Table~\ref{tab:joint} were first evaluated on the completed
experiments only.  Before the lost outcomes were known, a sensitivity showed that
every decision could change under a realizable assignment of them: if every
interrupted experiment had exceeded the observed statistic, no test would have
been rejected.  All lost pseudo-experiments were then rerun under a procedure that
was fixed and independently reviewed beforehand, and that required a rerun of
completed experiments to reproduce them exactly (16 of 16 did).  None of the
recovered experiments reaches the observed statistic, in any test or declared
variant, so every submitted pseudo-experiment has now been observed and no
assignment of the lost outcomes changes any decision.  The closest approach is in
a report-only robustness variant of the GENIE~+~MEC shape test, \num{4.6} below the
observed statistic, against an observed value of \num{868.6}, and the next is
\num{49.8} (GENIE CV shape, also a robustness variant).  With complete batches, the
sequential rule would have stopped the Tune~v1, GENIE CV, GENIE~+~MEC and GiBUU
nulls one batch earlier, at $B=1200$, with the same decisions.  This resolution is
report-only: the decisions remain those of the frozen evaluation.  An independent
recomputation confirms it against the frozen variant shifts.  The reading with the
recovered experiments added to the calibration was reproduced only by a standalone
replay of the same rules from inputs extracted with the frozen code, a consistency
check rather than a fully independent recomputation, and the earlier stopping points
rest on the campaign's self-validated re-application of the sequential rule.
% Source: RECORD-20261005-s5p-joint-5d-inference-result.md sec. 5;
% missing-sensitivity.json f48e16ef...; PROCEDURE-20261005-s5p-lost-seed-recovery.md;
% RECORD-20261006-s5p-lost-seed-recovery-resolution.md (b354cdf7); cross-check 00009056;
% Joseph's disposition DECISION-20261005-...-lost-seed-recovery.md sec. 5 (538739ff): "Certified,
% margin disclosed", required disclosures 1-5 all stated above.

\subsection{Detector-response condition}
\label{sec:response}

MINERvA's recoil-energy analyses assign a hadronic-energy response uncertainty
derived from test-beam calorimetry~\cite{MINERvA:2015testbeam}, separate from the
GEANT hadron-interaction reweights~\cite{MINERvA:2021owq,MINERvA:2022incl}.  The
declared nuisance model here includes the reweights but no calorimetric response
band.  Such a band would act on all five nulls together, through the shared
detector simulation and unfolding, and could move each statistic in either
direction.  The published comparison in muon kinematics alone, whose
observables do not use the reconstructed recoil energy, also disfavours
the same generator families~\cite{MINERvA:2021owq}, so the condition bears
chiefly on the shape information along the hadronic axes.  As
an exploratory sensitivity, not a validated uncertainty prescription, the data
are re-unfolded with the simulation's reconstructed recoil energy scaled by
$\pm\SI{4}{\percent}$, the level of data--simulation agreement reported for the
test-beam calorimetric response~\cite{MINERvA:2015testbeam}, and the frozen
tests are re-evaluated.
A re-unfold with unscaled recoil reproduces the frozen decisions,
with every claim $p$ value inside the range spanned by the data's rounding
variations.  At both $-4\%$ and $+4\%$ every one of the ten rejections holds: the
observed statistics move by up to about \SI{11}{\percent} but remain far above
their null distributions, and only NuWro's shape test changes its count, from
$k=1$ to $k=3$ of $1751$ at $+4\%$ ($p=0.0023$, against a Holm threshold of
$0.05$).  The variation is applied to the data side after the energy calibration, as a single
coherent scale on the reconstructed available energy and recoil energy without
per-particle response, and not to the null calibration, whose response to it is
unmeasured: the rejections are not removed by this data-side variation, which does
not establish that the response model is complete.  The
evaluation was reproduced independently.
% Source: publication/w2/W2B-REPORT-20261006.md (study/w2-recoil-response-20261006, 940d84aa);
% independent check REPORT-20261006-s5p-recompute-w2-recoil-check.md (s5p-parallel-recompute-20260928, 9ffc4ddd):
% AGREE 1379/1379 on all three copies, control PASS (fallback), 150/150 per-sign fields.

\section{Uncertainty status}
\label{sec:uncertainty-status}

A scalar five-dimensional covariance is adopted under a digest-bound exception,
and its \zprojCells-cell $(\Eavail,W)$ projection is built; no superseded or
historical covariance is used here, and the joint tests do not use it.  We quote
no significance from it, and the reason is measured rather than precautionary.
In a dedicated campaign of two companion covariances built at two
estimator-seed settings (products distinct from the adopted bytes), varying the
seed moves the projected uncertainty by \SI{\sprojMeasured}{\percent}, against a
\SI{\sprojBound}{\percent} movement limit fixed before production.  That movement
did not fall with ensemble size over the range tested:
\SI{\seedEffectNforty}{\percent} averaged over four disjoint 40-throw subsets,
\SI{\seedEffectNeighty}{\percent} over the two 80-throw halves, and
\SI{\sprojMeasured}{\percent} on the full 160-throw ensemble, fitted exponent
$\seedEffectExponent$.  A same-seed resampling floor, measured at $N=40$ and
$N=80$, falls steeply over that step --- \SI{\floorNforty}{\percent} to
\SI{\floorNeighty}{\percent}, two-point exponent $\floorExponent$.  The comparison
is a single seed pair and the 40- and 80-throw points are nested subsets of the
same 160 throws, so the width of the seed-pair distribution and the behaviour at
much larger $N$ are both unmeasured; the statistic is a maximum over the declared
offset set, so further pairs can only raise it.  Five of the covariance's 45
bands are built at a fixed seed the variation cannot reach, carrying
\SI{\pinnedBandsSqrtTr}{\percent} of $\sqrt{\mathrm{Tr}\,C}$ and contributing none of
the movement, so their seed sensitivity is unprobed and the effect on the total
were they also to vary is not measured in either direction.  The same variation
moves the projected central values by at most
\SI{\cvMoveOverSigmaMax}{\percent} of their own quoted uncertainty (median
\SI{\cvMoveOverSigmaMed}{\percent}).

A separate campaign designed to construct a five-dimensional measurement with
qualified total uncertainties did not admit any reporting definition: none could
meet its predeclared useful-precision targets with any configuration its rules
allowed, a scientific rather than a computational limitation.  The joint tests
above were its remaining,
separately admitted branch.  For the fallback definition, the
bounded model-dependence component alone exceeded the total-width target.  Full-event point-cloud (PET) classifiers that
could replace the scalar learner remain diagnostic and method development: the
historical nominal/statistical pairing was declined, no PET uncertainty product
is adopted, a subsequent improvement study adopted no candidate, and a
final-design study found no eligible design~\cite{PETImprovement2026,PETFinalDesign2026}.  An exploratory,
exact-paired simulation comparison with a HistGradientBoosting OmniFold (distinct
from the production LightGBM) found better recovery of an injected energy tilt
for PET than for that tree method (a binned unfolding of the same bins matched PET
on the tilt) but better recovery of all three tested generator reweightings for the
tree method~\cite{PETGBDT2026}; it establishes neither a general method ranking
nor an uncertainty comparison.
% Condensed from main's paper_body.tex at 30942b49 (PR #20, e55f6e9f); PET sources: VL164-VL167
% (final design, NO_ELIGIBLE_DESIGN), VL168 (PET vs GBDT), improvement REPORT-20260922.md sec. 12;
% the companion note carries the full PET results (sec_pet_campaigns.tex). VL142-VL144 in the
% comment above back the covariance sentences only.
% Sources: VL142-VL144; DECISION-20260920-joseph-adopts-z-cv-under-the-6.4-exception.md sec. 4;
% CORRECTION-20260920/-20260921; RECORD-20260927-s5p-stage2-exit.md; amendment 4;
% DECISION_RECORD-pet-final-design.md.

\section{Conclusions}
\label{sec:conclusions}

An unbinned event-weight unfolding reproduces the published two-dimensional
MINERvA inclusive measurement and extends it to five simultaneous observables
whose central values pass marginal anchors and closure tests and recover the
known low-recoil behavior.  The measured limitations --- a fixed-truth
undercoverage of the two-dimensional statistical band as first built (the rebuilt
band has not been re-tested), a background bias at
high $W$, and a regularization bias under generator-sized shape changes ---
are reported with the results they affect, and they prevent a qualified
five-dimensional measurement.  A joint test calibrated through the full chain
provides a statement about specified predictions that does not require such a
measurement.
Frozen joint tests of five neutrino-interaction predictions against the
five-dimensional data reject all ten simple hypotheses at familywise level 0.05
under the stated conditions.  Every pseudo-experiment first lost to scheduler
limits has been observed in a report-only recovery without changing a decision, an exploratory
$\pm4\%$ recoil-response variation leaves every rejection in place, and each
rejected test already has $p<0.05$ in at least one matched coarse projection, so
no claim is made that the joint information adds discrimination.  These are statements about the
specified predictions, not a measured cross section, and they add no new
reportable uncertainty.
A complete five-dimensional measurement release would additionally require a
matched total uncertainty with validated coverage.

\section{Data availability}
\label{sec:data-availability}

The input data and simulation are the MINERvA open-data AnaTuples, released under
CC0~\cite{MINERvA:OpenData}; this analysis used an earlier production of that
release whose file identities are recorded with the release.
A release of about \SI{16}{MB} lets a reader, with standard numerical libraries
only, recompute every $p$ value, decision, robustness label and power figure of the
joint tests, the lost-seed resolution and the coarse-projection comparison;
regenerate Figs.~\ref{fig:validation}--\ref{fig:joint-nulls}; and recompute the
numbers quoted from Figs.~\ref{fig:validation}--\ref{fig:generator-context}.  It
contains the per-pseudo-experiment cell values after the declared nuisance terms,
their seeds, the predictions, the test domains, the variant shifts, the metric and
the observed data vector, the arrays behind Figs.~\ref{fig:validation}--\ref{fig:generator-context},
and the digest of every source product.  Replays of the packaged files reproduced
every recorded value and every checked figure number, one of them by an outside
reader from an empty directory.  The remaining quoted numbers (the
two-dimensional uncertainty and coverage figures, the closure and bias studies,
the low-recoil comparisons, the recoil-response sensitivity and the covariance-status
measurements) are not recomputable from this release.
Extraction from the pseudo-experiment products, which are archived
(\SI{2.3}{GB}), and the event-level processing are not reproducible from it
either; they need the archived products and the open-data AnaTuples together
with the generator samples produced for this analysis.
\pubhold{deposit identifier and access route, after the separately authorized
release decision.}
% Source: docs/publication/release/RC3-README.md (RC3 7999d094..., VERIFY PASS); RC1-REPLAY-RECEIPT-20261006.md
% (outside empty-directory replay); final review cycle 2 re-ran RC3 verify (PASS); product archive
% /global/cfs/cdirs/m3246/josephrb/s5p-archive-20261006/ (SHA256SUMS 1a72cb43..., main 275b996d).
